\section{Unrestricted rank for block-exchangeable densities}
\label{sec:exchangeable}
Consider two copies of the identity column basis, with no other nonzero
columns. Half noise induces independent bit noise of correlation $\rho=1/4$
on the syndrome cube. Partition its $r$ coordinates into a fixed number $b$
of blocks of sizes $n_i\ge\gamma r$, with fixed $\gamma>0$. Suppose the
positive stationary density $\pi$ is invariant under permutations within
each block, $\|\pi\|_{15/7}\le B$, and stationary ambient marginals match
through order $k\to\infty$. There is no relation between $r$ and $k$.

\paragraph{Theorem.} For every sequence of syndrome sets $E$ with
$\mu(E)\to0$, the two compressed norms for $W=\pi/g^2$ and $W=\pi/V^2$
in the proof of (SS3) tend to zero. The sets and vectors in the compression
need not have any permutation symmetry. Consequently the full ordered
spectrum recovers that of ordinary noise, and the normalized trial has
vanishing return ratio on every subset of the critical set $D$.
This includes arbitrary radial densities ($b=1$), with no positive floor
or inverse entropy assumption. We prove qualitative convergence.

\paragraph{A lower bound for the transition to a central layer.}
In a block of size $n$, put $z=n^{-1/2}\sum x_j$ using signs $\pm1$.
Write $p_z(w)$ for the transition probability to the layer with
standardized sum $w$, and $p_*(w)$ for its uniform mass. Put
$q=1-\rho^2=15/16$. Fix $K$ and a finite $A$. For $|z|\le K n^{1/3}$
and attainable $|w|\le A+1$, for sufficiently large $n$,
\[
 \log \frac{p_z(w)}{p_*(w)}
 \ge -\frac{(w-\rho z)^2}{2q}
        -C_\rho\frac{(|z|+|w|)^4}{n}-C_\rho.                 \tag{BX1}
\]
For the unstandardized output sum $Y$, its log moment generating
function $\Lambda(t)$ is the sum of the logs of
$\cosh t+\rho x_j\sinh t$. Its first three derivatives at zero are
$\rho z\sqrt n$, $nq$, and $-2\rho q z\sqrt n$; its fourth derivative is
bounded by $C_\rho n$ near zero. Tilt so that $\Lambda'(t)=w\sqrt n$.
Since $\Lambda''\ge c_\rho n$ near zero, the solution satisfies
$|t|\le C_\rho(|z|+|w|)/\sqrt n$. Taylor's theorem and the maximum of
the quadratic in $t$ bound $tw\sqrt n-\Lambda(t)$ above by the first
two costs on the right of (BX1), taken with positive signs.
Under the tilt the numbers of output plus signs from the two input
sign groups are independent binomials. Both group sizes are comparable
to $n$ and both success probabilities stay bounded away from zero and one.
There are at least $c\sqrt n$ integer splits of the target count with
each count within $\sqrt n$ of its tilted mean: the two means sum to
the target. Stirling's bounds give mass at least $c/\sqrt n$ to each
such count in its group, hence at least $c/\sqrt n$ to their sum at
the target. Undo the tilt and use $p_*(w)\le C/\sqrt n$.
This elementary estimate uses the correlated-binomial representation;
see \href{https://doi.org/10.1111/j.1467-842X.2012.00654.x}{Diaconis
and Griffiths} for the associated bivariate binomial laws.

\paragraph{Matching supplies central layers on both sides.}
Write $Z=(Z_1,\ldots,Z_b)$ for the standardized sums. Under $\nu=\pi\mu$
the expectation of every fixed polynomial in $Z$ eventually equals its
uniform expectation: a syndrome character of degree $j$ has ambient
degree $2j$. For $r\to\infty$ the uniform mixed moments converge to
those of independent standard Gaussians by the product moment calculation.
The Gaussian is moment determinate (equivalently apply the scalar moment
method to each linear combination). Every fixed box with nonempty interior
and Gaussian-null boundary therefore has $\nu$-mass bounded below.
For each choice of signs take the box with signed coordinates in $[A,A+1]$.
Their number is fixed. The integral of $\sqrt\pi$ over each box, with
respect to $\mu$, is bounded below by row-free H\"older interpolation of
the first moment between $1/2$ and $15/7$. The same assertion for $h$
follows from $\|h-\pi\|_1\to0$ proved before (M3).
The positive scaling $h$ is block invariant by uniqueness and equivariance.

\paragraph{Endpoint cancellation.}
Fix $K$. Call the region $|Z_i|\le K n_i^{1/3}$ for every $i$ moderate.
For $f=h$ or $f=\sqrt\pi$, the block symmetry expresses $Pf(x)$ as the
integral of $f$ over layers against the product of the $b$ row-to-uniform
ratios. Restrict to the box whose signs agree with $z=Z(x)$.
In (BX1) the cross term is at least $\rho A |z_i|/q$, while
$|z_i|^4/n_i\le K^3|z_i|$. The coefficient of the quartic cost is
independent of $A$. Taking $A$ sufficiently large \emph{after} fixing
$K$, and absorbing terms depending only on $A,b,K,\rho$ in a constant,
gives uniformly on the moderate region, for all sufficiently large ranks,
\[
        (Pf(x))^{-15}\le C
        \exp\!\left(\tfrac12\sum_i z_i^2-3\sum_i|z_i|\right).
                                                               \tag{BX2}
\]
Indeed $15\rho^2/(2q)=1/2$ exactly. On the moderate region these powers
are uniformly integrable under $\mu$. To see this, use independent blocks
and $\mu(|Z_i|\ge m)\le 2e^{-m^2/2}$. The integral bound on
$m\le|Z_i|<m+1$ in each factor is at most $2\exp(1/2-2m)$.
The sum outside any fixed box tends uniformly to zero as the box grows;
inside the box (BX2) is bounded. Thus for every $\mu(E)\to0$,
$\|(Pf)^{-1}1_{E\cap{\rm moderate}}\|_{15}\to0$.

\paragraph{The outer layers and the full operator.}
The complement of the moderate region has uniform measure at most
$2b\exp[-K^2(\gamma r)^{2/3}/2]$. In (SS1) use $d=r$. Since
$L^2/(r+L)$ increases with $L$, choosing $K$ sufficiently large
depending on $B,b,\gamma$ makes the negative term dominate
$C_B r^{1/3}$ by a fixed positive multiple of $r^{1/3}$.
The operator-norm proof of (SS1) applies to both weights on the entire
outer region. On the moderate part of $E$ use
$\|1_{E\cap{\rm moderate}}W\|_{5/3}\le B
\|1_{E\cap{\rm moderate}}/(Pf)\|_{15}^2$ and the contractive
$2\to5$ factors of (CR3). For their union the norm is at most the sum
of the two bounds: the nonzero spectrum of the compression equals that
of $P^{1/2}M_{W1_E}P^{1/2}$, which is additive in the restricted weight.
The trial normalization is bounded as in (SS1). The set $D$ has
$\mu(D)\to0$ by its high-density threshold. The high-multiplier set of
(FF1) also has vanishing uniform measure by its fixed positive moment
bound before (CR1). Applying (FF1) proves full ordered recovery, including
the ambient fiber modes covered by that criterion. On bounded-rank
subsequences matching eventually implies $\pi=1$ and recovery is exact.
These arguments use balanced block sizes and a product syndrome kernel;
a general pair of bases need not have this layer transition law.
